% ============================================================
\chapter{Introduction}
\label{chap:intro}
% ============================================================

\section{Problem Statement}

A developer with a laptop, a laboratory workstation and a spare desktop has a
considerable amount of idle computing capacity and no practical way to use it.
The frameworks that exist for distributed computation --- Apache Spark
\cite{zaharia2010spark}, Ray \cite{moritz2018ray}, Dask \cite{rocklin2015dask}
--- presuppose a cluster that somebody has already configured: a head node,
matching software on every worker, a scheduler process, and a network on which
every machine can reach every other. Ordinary machines fail the last of these,
since each sits behind its own home or campus router and accepts no inbound
connection. The alternative, \texttt{ssh} with \texttt{rsync}, means reproducing
the environment by hand on every machine and receiving no scheduling, isolation
or fault tolerance in return.

Three distinct problems sit underneath this:

\begin{enumerate}[leftmargin=*, label=\arabic*.]
  \item \textbf{The environment problem.} A script needs a particular
  interpreter and a particular set of libraries, and making a remote machine
  match the local one is the bulk of the work.
  \item \textbf{The code problem.} Even with a matching environment, using more
  than one machine normally means rewriting the program against a cluster API.
  \item \textbf{The reachability problem.} Machines on different networks cannot
  open a connection to one another without a VPN, port forwarding, or a relay
  that carries every byte between them.
\end{enumerate}

\textbf{PipedPeer} addresses all three without asking the user to change the
program. The environment is reproduced exactly by shipping a Nix closure; the
program is distributed without modification by intercepting the parallel
primitives it already uses; and machines on different networks are joined by a
direct path through their routers, with an introducer that only tells each
machine where the others are. No machine is a head node: coordination is done
per job by the machine that submits it.

\section{Objectives}

\begin{enumerate}[leftmargin=*, label=\arabic*.]
  \item To run an unmodified Python script on the best available machine, with
  results appearing in the user's working directory as though the script had
  run locally.
  \item To reproduce the script's environment exactly on that machine, sending
  only the store paths that machine does not already hold.
  \item To choose the machine from live capacity, treating architecture and GPU
  availability as hard requirements and load, memory and hardware as a ranking.
  \item To distribute the parallel work inside the script with no code changes,
  only when doing so is expected to be faster than running it locally, and in
  proportion to each machine's measured throughput.
  \item To join machines on different networks into one cluster with a single
  command, with no VPN, no port forwarding and no relay in the data path.
  \item To authenticate every request, encrypt traffic between machines, isolate
  each job, and survive the loss of any worker without losing the run.
  \item To require no manual installation on a new machine: one binary, whose
  set-up command installs what else is needed, without root if necessary.
\end{enumerate}

\section{Scope}

Within scope and implemented are a Go daemon and CLI in one binary; discovery
on the local subnet and an address book of every path to every node; joining
across networks by an introducer in every daemon, router port mapping, hole
punching and QUIC links; a shared secret, pinned TLS and signed closures;
reproducible Nix environments transferred as a path-level difference; placement
by hard filters and a score, with a three-phase lease; per-job \texttt{crun}
sandboxes under memory ceilings, with GPU passthrough and without root;
transparent interception of \texttt{multiprocessing.Pool}, pandas, NumPy,
\texttt{joblib} and scikit-learn, and torch data-parallel training; and Android
phones as workers for parallel chunks. Out of scope are per-user authorisation (a cluster shares one secret;
Section~\ref{sec:security}), runtimes other than Python, relaying traffic for a
peer that no direct path can reach (such a peer is reported and routed around),
and any incentive mechanism for contributing capacity.

\section{Technologies Used}

\begin{table}[!htbp]
\centering
\caption{Technologies used in PipedPeer}
\label{tab:tech}
\renewcommand{\arraystretch}{1.2}
\rowcolors{2}{}{rowgray}
\begin{tabularx}{\textwidth}{L{4.3cm} L{2.4cm} X}
\toprule
\textbf{Technology} & \textbf{Category} & \textbf{Purpose} \\
\midrule
Go 1.25 & Core language & Daemon, CLI and coordinator in one static binary \\
\texttt{go-chi/chi} & HTTP & Routing for the daemon API \\
\texttt{gorilla/websocket} & Streaming & Live output from a running job \\
\texttt{quic-go} & Transport & Encrypted, multiplexed peer links over the punched UDP socket \\
\texttt{crypto/tls}, Ed25519 & Security & Pinned TLS on the daemon port; node identity and signing keys \\
Nix (pinned) & Environments & Content-addressed closures; the store path is the cache key \\
\texttt{crun} & Sandboxing & OCI container per job, including GPU device passthrough \\
cgroup v2, \texttt{systemd-run} & Resources & Memory ceilings for the daemon and for each job \\
zstd (\texttt{klauspost/compress}) & Compression & Workspace, results and request bodies \\
\texttt{modernc.org/sqlite} & Storage & Address book, in pure Go so the binary stays static \\
\texttt{gopsutil} & Metrics & CPU, memory and process-tree measurements \\
Cobra; Bubble Tea, Lip Gloss & CLI, TUI & Command structure; live dashboard \\
\texttt{zerolog} & Logging & Structured, levelled logging \\
Python, cloudpickle & Interception & The shim, embedded in the binary; shipping closures by value \\
\bottomrule
\end{tabularx}
\end{table}

\section{Assumptions and Constraints}

\begin{itemize}[leftmargin=*]
  \item Every participating machine runs the same binary.
  \texttt{pipedpeer setup} installs Nix and \texttt{crun} when they are missing,
  and can do so without root by using a private Nix store.
  \item Workers run Linux, since \texttt{crun}, cgroups and the Nix store are
  Linux-specific; phones run Linux user space under Termux.
  \item Machines in a cluster share one secret, created by the first machine's
  set-up and carried to the others by an invite code.
  \item Joining across networks needs one reachable member or public
  introducer. Introducers carry no job data.
  \item Two machines between which no direct path exists cannot work together.
  \item The submitting process holds the lease of a parallel script and must
  remain alive for its duration; a single remote job can be detached.
\end{itemize}

\section{Layout of the Report}

\begin{itemize}[leftmargin=*]
  \item Chapter~\ref{chap:lit} surveys the work PipedPeer draws on.
  \item Chapter~\ref{chap:arch} gives the design: architecture, data flow, use
  cases and the sequence of a run.
  \item Chapter~\ref{chap:design} describes the implementation of each
  mechanism, from discovery and joining across networks to interception,
  training and fault tolerance.
  \item Chapter~\ref{chap:results} sets out the evaluation and its results.
  \item Chapter~\ref{chap:conclusion} concludes, with limitations and security.
  \item Chapter~\ref{sec:future} lists future scope, and
  Chapter~\ref{chap:pub} summarises the research paper drawn from this work.
  \item Appendix~\ref{app:extra} holds the implementation's size and further
  measurements.
\end{itemize}
